\documentclass[10pt,a4paper]{amsart}
\usepackage{amsmath,amssymb,amsthm,graphicx}
\usepackage{verbatim}
\newenvironment{reference}{\comment}{\endcomment}
%\newenvironment{reference}{}{}
\newenvironment{slogan}{\comment}{\endcomment}
\newenvironment{history}{\comment}{\endcomment}
% For commutative diagrams we use Xy-pic
\usepackage[all]{xy}

% We use 2cell for 2-commutative diagrams.
\xyoption{2cell}
% Theorem environments.
%
\theoremstyle{plain}
\newtheorem{theorem}[subsection]{Theorem}
\newtheorem{proposition}[subsection]{Proposition}
\newtheorem{lemma}[subsection]{Lemma}

\theoremstyle{definition}
\newtheorem{definition}[subsection]{Definition}
\newtheorem{example}[subsection]{Example}
\newtheorem{exercise}[subsection]{Exercise}
\newtheorem{situation}[subsection]{Situation}

\theoremstyle{remark}
\newtheorem{remark}[subsection]{Remark}
\newtheorem{remarks}[subsection]{Remarks}

\numberwithin{equation}{subsection}

% Macros
%
\def\lim{\mathop{\mathrm{lim}}\nolimits}
\def\colim{\mathop{\mathrm{colim}}\nolimits}
\def\Spec{\mathop{\mathrm{Spec}}}
\def\Hom{\mathop{\mathrm{Hom}}\nolimits}
\def\Ext{\mathop{\mathrm{Ext}}\nolimits}
\def\SheafHom{\mathop{\mathcal{H}\!\mathit{om}}\nolimits}
\def\SheafExt{\mathop{\mathcal{E}\!\mathit{xt}}\nolimits}
\def\Sch{\mathit{Sch}}
\def\Mor{\mathop{\mathrm{Mor}}\nolimits}
\def\Ob{\mathop{\mathrm{Ob}}\nolimits}
\def\Sh{\mathop{\mathit{Sh}}\nolimits}
\def\NL{\mathop{N\!L}\nolimits}
\def\CH{\mathop{\mathrm{CH}}\nolimits}
\def\proetale{{pro\text{-}\acute{e}tale}}
\def\etale{{\acute{e}tale}}
\def\QCoh{\mathit{QCoh}}
\def\Ker{\mathop{\mathrm{Ker}}}
\def\Im{\mathop{\mathrm{Im}}}
\def\Coker{\mathop{\mathrm{Coker}}}
\def\Coim{\mathop{\mathrm{Coim}}}

% Boxtimes
%
\DeclareMathSymbol{\boxtimes}{\mathbin}{AMSa}{"02}

%
% Macros for moduli stacks/spaces
%
\def\QCohstack{\mathcal{QC}\!\mathit{oh}}
\def\Cohstack{\mathcal{C}\!\mathit{oh}}
\def\Spacesstack{\mathcal{S}\!\mathit{paces}}
\def\Quotfunctor{\mathrm{Quot}}
\def\Hilbfunctor{\mathrm{Hilb}}
\def\Curvesstack{\mathcal{C}\!\mathit{urves}}
\def\Polarizedstack{\mathcal{P}\!\mathit{olarized}}
\def\Complexesstack{\mathcal{C}\!\mathit{omplexes}}
% \Pic is the operator that assigns to X its picard group, usage \Pic(X)
% \Picardstack_{X/B} denotes the Picard stack of X over B
% \Picardfunctor_{X/B} denotes the Picard functor of X over B
\def\Pic{\mathop{\mathrm{Pic}}\nolimits}
\def\Picardstack{\mathcal{P}\!\mathit{ic}}
\def\Picardfunctor{\mathrm{Pic}}
\def\Deformationcategory{\mathcal{D}\!\mathit{ef}}
\usepackage{cleveref}
\crefname{originalreference}{original reference}{original references}
\Crefname{originalreference}{Original reference}{Original references}
\makeatletter
\newlabel{algebra-lemma-colimit-syntomic}{{algebra-lemma-colimit-syntomic}{source}{}{sourcepointer-1}{}}
\newlabel{algebra-lemma-colimit-syntomic@cref}{{[originalreference][0][]algebra-lemma-colimit-syntomic}{[1][0][]source}{}{}{}}
\makeatother
\crefname{originalreference}{original reference}{original references}
\Crefname{originalreference}{Original reference}{Original references}
\makeatletter
\newlabel{dpa-definition-lci}{{dpa-definition-lci}{source}{}{sourcepointer-2}{}}
\newlabel{dpa-definition-lci@cref}{{[originalreference][0][]dpa-definition-lci}{[1][0][]source}{}{}{}}
\makeatother
\crefname{originalreference}{original reference}{original references}
\Crefname{originalreference}{Original reference}{Original references}
\makeatletter
\newlabel{dpa-lemma-avramov}{{dpa-lemma-avramov}{source}{}{sourcepointer-3}{}}
\newlabel{dpa-lemma-avramov@cref}{{[originalreference][0][]dpa-lemma-avramov}{[1][0][]source}{}{}{}}
\makeatother
\crefname{originalreference}{original reference}{original references}
\Crefname{originalreference}{Original reference}{Original references}
\makeatletter
\newlabel{more-algebra-section-properties-formal-fibres}{{more-algebra-section-properties-formal-fibres}{source}{}{sourcepointer-4}{}}
\newlabel{more-algebra-section-properties-formal-fibres@cref}{{[originalreference][0][]more-algebra-section-properties-formal-fibres}{[1][0][]source}{}{}{}}
\makeatother
\title{Being a local complete intersection is a property of the formal fibres}
\author{The Stacks Project Authors}
\date{stacks-project:dualizing.tex:local-snapshot-20261012; GNU FDL 1.2 or later}
\begin{document}
\maketitle
\noindent\textit{Editorial note.} All delivered mathematical text is the author's own, copied from the pinned Stacks Project chapter source sources/stacks-dualizing/pinned-original-source.tex; no statement or proof step was rewritten, reordered or removed.\par\medskip
\noindent\textit{Editorial note.} Preamble provenance: the Stacks Project keeps one single shared preamble for every chapter; this card takes it directly from the staged sources/stacks-dualizing/preamble.tex as three verbatim spans (lines 8-12, 14-18 and 154-217), so no derived file was written and nothing was added or redefined. The wrapper class and the base packages are supplied by the assembler.\par\medskip
\noindent\textit{Editorial note.} Bibliography provenance: neither the counted statement nor its proof contains a citation (this lane rescans the whole span before writing the card), so the card carries no bibliography and no entry was reconstructed.\par\medskip
\noindent\textit{Editorial note.} Selection provenance: this card was opened under the file-based control contract (CONTROL.md), with the mechanical pre-card gate run on its statement and proof interval before the card was written.\par\medskip
\noindent\textit{Allowed standard prerequisites (editorial).}\begin{itemize}
\item The presentation of a finitely generated field extension as the fraction field of a global complete intersection. Invoked inside the counted proof (source lines 4498-4501). Supplies the auxiliary algebra over which the base change is a relative global complete intersection.
\item The permanence of the complete intersection property under a relative global complete intersection. Invoked inside the counted proof (source lines 4502-4504 and 4519-4520). Carries the property through the base change and through the regular extension in part (C).
\end{itemize}\medskip
\ifdefined\hypertarget\hypertarget{sourcepointer-1}{}\fi
\noindent\textit{Source reference.} Original reference algebra-lemma-colimit-syntomic: sources/stacks-dualizing/pinned-original-source.tex, source lines 4426--4426. Navigation pointer only: the counted text refers to the Stacks item algebra-lemma-colimit-syntomic, which belongs to another chapter of the project and is not part of this card. The pointer carries the target's own label as its number because this chapter's staged snapshot does not contain the other chapter's tag table; its content is not reproduced here.\par\medskip
\ifdefined\hypertarget\hypertarget{sourcepointer-2}{}\fi
\noindent\textit{Source reference.} Original reference dpa-definition-lci: sources/stacks-dualizing/pinned-original-source.tex, source lines 4492--4492. Navigation pointer only: the counted text refers to the Stacks item dpa-definition-lci, which belongs to another chapter of the project and is not part of this card. The pointer carries the target's own label as its number because this chapter's staged snapshot does not contain the other chapter's tag table; its content is not reproduced here.\par\medskip
\ifdefined\hypertarget\hypertarget{sourcepointer-3}{}\fi
\noindent\textit{Source reference.} Original reference dpa-lemma-avramov: sources/stacks-dualizing/pinned-original-source.tex, source lines 4504--4504. Navigation pointer only: the counted text refers to the Stacks item dpa-lemma-avramov, which belongs to another chapter of the project and is not part of this card. The pointer carries the target's own label as its number because this chapter's staged snapshot does not contain the other chapter's tag table; its content is not reproduced here.\par\medskip
\ifdefined\hypertarget\hypertarget{sourcepointer-4}{}\fi
\noindent\textit{Source reference.} Original reference more-algebra-section-properties-formal-fibres: sources/stacks-dualizing/pinned-original-source.tex, source lines 4385--4385. Navigation pointer only: the counted text refers to the Stacks item more-algebra-section-properties-formal-fibres, which belongs to another chapter of the project and is not part of this card. The pointer carries the target's own label as its number because this chapter's staged snapshot does not contain the other chapter's tag table; its content is not reproduced here.\par\medskip
\begin{lemma}
\label{lemma-formal-fibres-lci}
Properties (A), (B), (C), (D), and (E) of
More on Algebra, Section \ref{more-algebra-section-properties-formal-fibres}
hold for $P(k \to R) =$``$R$ is a local complete intersection''.
See Divided Power Algebra, Definition \ref{dpa-definition-lci}.
\end{lemma}

\begin{proof}
Part (A). Let $K/k$ be a finitely generated field extension.
Let $R$ be a $k$-algebra which is a local complete intersection.
We can find a global complete intersection
$A = k[x_1, \ldots, x_n]/(f_1, \ldots, f_c)$
over $k$ such that $K$ is isomorphic to the fraction field of $A$, see
Algebra, Lemma \ref{algebra-lemma-colimit-syntomic}.
Then $R \to R \otimes_k A$ is a relative global complete intersection.
It follows that $R \otimes_k A$ is a local complete intersection
by Divided Power Algebra, Lemma \ref{dpa-lemma-avramov}.

\medskip\noindent
Proof of (B). This is clear
because a ring is a local complete intersection if and only if all of its
local rings are complete intersections.

\medskip\noindent
Part (C). Let $A \to B \to C$ be flat maps of Noetherian rings.
Assume the fibres of $A \to B$ are local complete intersections
and $B \to C$ is regular. We have to show the fibres of $A \to C$
are local complete intersections. Clearly, we may assume $A = k$ is a field.
Then we may assume that $B \to C$ is a regular local homomorphism
of Noetherian local rings. Then $B$ is a complete intersection and
$C/\mathfrak m_B C$ is regular, in particular a complete intersection
(by definition). Then $C$ is a complete intersection by
Divided Power Algebra, Lemma \ref{dpa-lemma-avramov}.

\medskip\noindent
Part (D). This follows by the same arguments as in (C) from
the other implication in
Divided Power Algebra, Lemma \ref{dpa-lemma-avramov}.
Part (E) is immediate as the condition does not refer to the ground
field.
\end{proof}

\end{document}
